\documentclass[11pt]{article}
\usepackage[margin=1in]{geometry}
\usepackage{times}
\usepackage{amsmath,amssymb}
\usepackage{graphicx}
\usepackage{booktabs}
\usepackage{hyperref}
\usepackage[numbers,sort&compress]{natbib}
\usepackage{authblk}

\hypersetup{colorlinks=true, linkcolor=blue, citecolor=blue, urlcolor=blue}

% Every empirical number below is generated from the committed result files by
% make_tables.py (see README). Set the repository URL once here.
\newcommand{\RepoURL}{https://github.com/USERNAME/REPO}
\newcommand{\resfile}[1]{\href{\RepoURL/blob/main/results/#1}{\nolinkurl{results/#1}}}
% generated by make_tables.py from results/*.json -- do not edit by hand
\newcommand{\TFullInstTrained}{0.994}
\newcommand{\TFullInstTrainedSd}{0.003}
\newcommand{\TFullInstRandom}{0.993}
\newcommand{\TFullInstRandomSd}{0.002}
\newcommand{\TFullInstGap}{+0.10}
\newcommand{\TFullInstGapSd}{0.15}
\newcommand{\TFullInstZone}{0.999}
\newcommand{\TFullInstLevelUse}{0.0203}
\newcommand{\TFullInstGapMin}{-0.1}
\newcommand{\TFullInstGapMax}{0.3}
\newcommand{\TFullInstN}{3}
\newcommand{\TFullRecTrained}{0.994}
\newcommand{\TFullRecTrainedSd}{0.002}
\newcommand{\TFullRecRandom}{0.985}
\newcommand{\TFullRecRandomSd}{0.003}
\newcommand{\TFullRecGap}{+0.83}
\newcommand{\TFullRecGapSd}{0.17}
\newcommand{\TFullRecZone}{0.999}
\newcommand{\TFullRecLevelUse}{0.0188}
\newcommand{\TFullRecGapMin}{0.6}
\newcommand{\TFullRecGapMax}{1.0}
\newcommand{\TFullRecN}{3}
\newcommand{\TEgoInstTrained}{0.401}
\newcommand{\TEgoInstTrainedSd}{0.043}
\newcommand{\TEgoInstRandom}{0.396}
\newcommand{\TEgoInstRandomSd}{0.060}
\newcommand{\TEgoInstGap}{+0.51}
\newcommand{\TEgoInstGapSd}{2.23}
\newcommand{\TEgoInstZone}{0.437}
\newcommand{\TEgoInstLevelUse}{0.0362}
\newcommand{\TEgoInstGapMin}{-1.1}
\newcommand{\TEgoInstGapMax}{3.7}
\newcommand{\TEgoInstN}{3}
\newcommand{\TEgoRecTrained}{0.654}
\newcommand{\TEgoRecTrainedSd}{0.060}
\newcommand{\TEgoRecRandom}{0.455}
\newcommand{\TEgoRecRandomSd}{0.067}
\newcommand{\TEgoRecGap}{+19.91}
\newcommand{\TEgoRecGapSd}{3.36}
\newcommand{\TEgoRecZone}{0.451}
\newcommand{\TEgoRecLevelUse}{0.0597}
\newcommand{\TEgoRecGapMin}{15.9}
\newcommand{\TEgoRecGapMax}{24.1}
\newcommand{\TEgoRecN}{3}
\newcommand{\TLevelUseRatio}{1.6}
\newcommand{\MEgoNoMem}{0.810}
\newcommand{\MEgoNoMemSd}{0.015}
\newcommand{\MEgoMem}{0.899}
\newcommand{\MEgoMemSd}{0.006}
\newcommand{\MFullNoMem}{0.996}
\newcommand{\MFullNoMemSd}{0.0004}
\newcommand{\MFullMem}{0.995}
\newcommand{\MFullMemSd}{0.002}
\newcommand{\LZeroMean}{5.77}
\newcommand{\LZeroStd}{2.45}
\newcommand{\LThreeMean}{11.97}
\newcommand{\LThreeStd}{4.49}
\newcommand{\LSixMean}{4.03}
\newcommand{\LSixStd}{2.10}
\newcommand{\LTenMean}{4.68}
\newcommand{\LTenStd}{2.13}
\newcommand{\LFifteenMean}{5.41}
\newcommand{\LFifteenStd}{5.35}
\newcommand{\LZeroSeedZero}{2.3}
\newcommand{\LThreeSeedZero}{5.7}
\newcommand{\LSixSeedZero}{6.3}
\newcommand{\EnvTransitions}{1.0}
\newcommand{\EnvEpLen}{75}
\newcommand{\EnvNoTransitionPct}{48}

\makeatletter\newcommand{\tablerows}[1]{\@@input #1 }\makeatother

\title{When Does Hierarchical Abstraction in World Models Actually Learn Anything?\\
\large A Controlled Diagnostic Study}

\author[1]{Regina Nam}
\affil[1]{Ionosphere Institute}
\date{}

\begin{document}
\maketitle

\begin{abstract}
Hierarchical predictive architectures such as H-JEPA \citep{lecun2022path} propose learning abstractions at multiple temporal and spatial scales, and have recently been instantiated end-to-end for visual planning \citep{zhang2026hjepa}. Evidence that such a hierarchy has learned an abstraction typically comes from probing: a slow variable remains decodable at higher levels while fast detail does not. Probe recoverability, however, does not distinguish information a level \emph{learned} from information it merely \emph{preserved} from the level below. We introduce a diagnostic protocol built around a \emph{random-abstractor control} --- comparing a trained top level against an identically-parameterized but untrained one --- and apply it across a $2\times2$ design crossing observability (full vs.\ egocentric) with abstractor type (instantaneous vs.\ recurrent) in controlled gridworld environments with exact ground truth at both scales. Three of four conditions yield abstractions that are statistically indistinguishable from random projections (or slightly worse), for two distinct reasons: under full observability the abstraction has nothing to add, since the slow variable is already linearly decodable from the fast level; under partial observability an instantaneous abstractor lacks any mechanism to integrate the history required to recover it. Only the combination of partial observability and a recurrent abstractor produces a measurably learned abstraction, replicated across \TEgoRecN{} seeds (top-level probe accuracy gap over random: $\TEgoRecGap\pm\TEgoRecGapSd$pp, vs.\ $\TFullRecGap\pm\TFullRecGapSd$pp under full observability). Current end-to-end hierarchical JEPAs, with pointwise level encoders in fixed-camera environments, sit in the cell where a level can reshape or discard information but cannot add it. We further report a negative result: attempts to establish a quantitative relationship between the density of persistent visual landmarks and abstraction quality in an open-world setting did not survive multi-seed replication, despite a promising single-run trend. Our results give an empirical account of \emph{why} architectures such as DreamerV3's RSSM \citep{hafner2023dreamerv3} combine a recurrent deterministic state with an instantaneous latent, and provide a reusable diagnostic that we argue should accompany any claim that a hierarchical model has learned a meaningful abstraction.
\end{abstract}

\section{Introduction}

Predicting in representation space rather than pixel space is the core commitment of the JEPA family of world models \citep{lecun2022path}: rather than reconstructing future observations, the model learns to predict future \emph{embeddings}, avoiding spending capacity on unpredictable detail. This principle has been demonstrated at scale for static images \citep{assran2023ijepa} and video \citep{bardes2024vjepa}. LeCun's 2022 position paper extends the idea to H-JEPA --- a hierarchical variant predicting at multiple temporal and spatial scales, with a slow abstract level providing context to a fast local one.

H-JEPA was introduced as a normative blueprint, without a concrete training loss for the abstract level. That gap has recently been closed: \citet{zhang2026hjepa} train a hierarchy of action-conditioned JEPAs end-to-end, each level predicting in its own latent space at a coarser temporal stride, and show that hierarchical planning outperforms flat planning on navigation and manipulation tasks. Their evidence that higher levels form abstractions comes from probing: across levels, slow state (the agent's position) remains decodable while fast detail (body pose, a moving distractor) becomes less so.

A training recipe does not settle a second question, which is the subject of this paper: how to \emph{check} that a level has learned something rather than passed it through. A probe that decodes a slow variable from a higher level establishes that the variable is recoverable there, not that training put it there --- a random projection of a lower level that already encodes the variable would score just as well. This matters because hierarchical models are prone to a specific failure: the lower level solves the task alone, gradients through the upper level become uninformative, and aggregate metrics still look healthy because the lower level carries them. Discarding fast detail, as observed by \citet{zhang2026hjepa}, is itself a training effect that a random projection would not be expected to reproduce; whether a level can \emph{add} information absent from the level below is a separate question, and the one our design isolates.

\paragraph{Contributions.}
\begin{enumerate}
\item A diagnostic protocol for hierarchical world models centered on a \emph{random-abstractor control}, plus two supporting diagnostics (level-use ablation and equal-capacity differential probing).
\item A $2\times2$ empirical study crossing observability with abstractor type, showing that three of four conditions produce abstractions statistically indistinguishable from random projections --- for two structurally different reasons --- and locating current end-to-end hierarchical JEPAs within this design.
\item An honest negative result on landmark density in open-world settings, including the methodological failure mode (a convincing trend on single runs that did not survive seed replication) that produced it.
\end{enumerate}

\section{Related Work}

\textbf{Self-supervised predictive representation learning.} VICReg \citep{bardes2021vicreg} regularizes joint-embedding representations via variance and covariance penalties to prevent collapse without negative pairs or a momentum target, and is the loss we build on throughout. I-JEPA \citep{assran2023ijepa} and V-JEPA \citep{bardes2024vjepa} demonstrate that predicting masked regions in representation space, rather than pixel space, yields semantic representations at scale for images and video respectively.

\textbf{Model-based RL with latent dynamics.} The Dreamer line of work \citep{hafner2023dreamerv3} learns a recurrent state-space model (RSSM) combining a deterministic recurrent state with a stochastic latent, and trains behavior entirely inside imagined rollouts. DreamerV3's architectural choice --- pairing recurrence with an instantaneous component --- is stated without an ablation isolating why recurrence is needed; our controlled study provides one candidate empirical explanation.

\textbf{Hierarchical abstraction.} H-JEPA \citep{lecun2022path} proposes hierarchical joint-embedding prediction as a path toward multi-timescale reasoning. \citet{zhang2026hjepa} provide an end-to-end instantiation: level encoders pool fixed windows of lower-level latents (pointwise, window size 1, in all reported experiments), each level is trained with a latent prediction loss and the SIGReg regularizer, and planning proceeds top-down with upper-level predictions serving as subgoals. They attribute planning gains to two mechanisms, a more informative goal-cost geometry in higher latent spaces and temporal decomposition into subgoal problems, and measure abstraction with nonlinear probes compared across levels. To our knowledge, no prior work, including \citet{zhang2026hjepa}, compares a trained hierarchical level against a random projection of the same architecture, or crosses abstractor recurrence with observability.

\section{Method}

\subsection{Environments}

All environments render $64\times64$ grayscale observations and expose ground truth at two scales, enabling exact probing of both levels.

\textbf{Building ($3\times3$ rooms).} A point agent moves in a grid of nine rooms connected by doorways. The \emph{fast} variable is position within the current room; the \emph{slow} variable is room identity, which changes rarely (on average \EnvTransitions{} room transitions per \EnvEpLen-step episode under the inertial random policy used for data collection, with no transition at all in \EnvNoTransitionPct\% of episodes; \resfile{env_stats.json}). Two observation modes: \texttt{full} (top-down view of the whole map) and \texttt{ego} (a window cropped around the agent, so rooms are mutually indistinguishable from a single frame and room identity is only recoverable by integrating history).

\textbf{Open world.} A larger ($128\times128$) continuous space scattered with circular obstacles and no room structure; the ``slow variable'' is an artificial $4\times4$ zone grid imposed only for diagnostic probing, not present in the environment's dynamics. Three variants: obstacles fixed across the dataset (landmark shortcuts possible), reshuffled every episode (pure dead reckoning), and reshuffled with $n$ persistent landmarks retained.

\subsection{Architecture}

The fast level is a convolutional encoder $z_1 = \text{Enc}(o)$ (128-d) with a residual dynamics predictor conditioned on both action and abstract context: $\hat z_{1,t+1} = z_{1,t} + f(z_{1,t}, a_t, z_{2,t})$. The abstract level is either

\begin{itemize}
\item \textbf{instantaneous}: $z_{2,t} = \text{MLP}(z_{1,t})$ (32-d), or
\item \textbf{recurrent}: $z_{2,t} = \text{GRU}(z_{1,t}, z_{2,t-1})$ (32-d),
\end{itemize}

with a jumpy level-2 predictor over $k$-step intervals taking an aggregated action summary as input. Both levels are trained with VICReg \citep{bardes2021vicreg} (invariance + variance + covariance) in latent space; no pixel reconstruction is used anywhere in training.

\subsection{Diagnostics}

\textbf{Random-abstractor control (primary).} We instantiate a second abstractor with identical architecture and dimensionality but untrained (random initialization), encode the same data through it, and fit the same probe. If trained $\approx$ random, training moved the abstraction nowhere a random projection of the fast level could not already reach. This control is the core of our protocol: without it, a high top-level probe score establishes only that the slow variable is \emph{recoverable}, not that it was \emph{learned}.

\textbf{Level-use ablation.} During training we measure the increase in level-1 prediction error when $z_2$ is replaced by another batch element's --- a hierarchical analogue of an action-conditioning ablation. Near-zero indicates the lower level ignores the abstract context.

\textbf{Equal-capacity differential probing.} Because $z_2 = g(z_1)$ for a nonlinear $g$, a \emph{linear} probe on $z_2$ is effectively a \emph{nonlinear} probe on $z_1$, and comparing the two directly is confounded. We therefore additionally probe both levels with matched-capacity nonlinear probes.

\section{Experiments}

\subsection{The $2\times2$ design}

Table~\ref{tab:2x2} reports linear-probe accuracy for the slow variable (room identity, 9 classes) on the fast level $z_1$, on the trained top level, and on the random control, for all four conditions over \TEgoRecN{} seeds (mean $\pm$ std).

\begin{table}[h]
\centering
\caption{Room-identity probe accuracy: fast level $z_1$, trained top level $z_2$, and random-abstractor control. Mean $\pm$ std over \TEgoRecN{} seeds. Source: \resfile{table1_2x2.json}.}
\label{tab:2x2}
\begin{tabular}{llcccc}
\toprule
Observability & Abstractor & $z_1$ & Trained $z_2$ & Random $z_2$ & Gap \\
\midrule
\tablerows{generated/table1.tex}
\bottomrule
\end{tabular}
\end{table}

Only one condition shows a learned abstraction. The three failures have two distinct causes.

\textbf{Nothing to add (full observability).} Room identity is already almost perfectly linearly decodable from $z_1$ (probe accuracy \TFullInstZone{} and \TFullRecZone{} in the two full-observability conditions), so any projection --- trained or random --- preserves it. The near-perfect scores in the top two rows measure \emph{recoverability}, not learning. This is the failure mode the random control was designed to catch, and without it these rows would read as strong successes.

\textbf{No mechanism (egocentric + instantaneous).} With an egocentric crop, room identity is not recoverable from any single frame; it requires integrating a history of door transitions. An instantaneous $z_2 = \text{MLP}(z_1)$ is a deterministic function of the current frame alone and therefore cannot encode more than $z_1$ does. The trained abstractor scores \emph{marginally below} random here, consistent with there being no learnable signal to acquire.

\textbf{Both conditions met (egocentric + recurrent).} Making the abstractor recurrent supplies the missing mechanism, and the $\TEgoRecGap\pm\TEgoRecGapSd$pp gap over random indicates a genuinely learned abstraction. The seed-to-seed range (\TEgoRecGapMin{} to \TEgoRecGapMax{}pp) is wide in absolute terms, underscoring why we replicate rather than report one number. The level-use ablation moves in the same direction (\TEgoInstLevelUse{} $\to$ \TEgoRecLevelUse{}, a $\TLevelUseRatio\times$ increase over the instantaneous case), indicating the lower level begins to rely on the abstract context once that context carries information the frame does not.

\subsection{Memory ablation with an observability control}

A separate experiment isolates memory in the \emph{dynamics} predictor (rather than the abstractor), in a two-room environment observed either top-down or through an egocentric crop. An identical pipeline is trained with a GRU-carrying versus a memoryless predictor, and global position is probed from $z$ for the memoryless model and from $(z, h)$ for the recurrent one (Table~\ref{tab:memory}).

\begin{table}[h]
\centering
\caption{Probe $R^2$ on global position: memoryless vs.\ recurrent predictor. Mean $\pm$ std over 3 seeds. Sources: \resfile{table2_memory_egocentric.json}, \resfile{table2_memory_base.json}.}
\label{tab:memory}
\begin{tabular}{lcc}
\toprule
Environment & No memory & With memory \\
\midrule
\tablerows{generated/table2.tex}
\bottomrule
\end{tabular}
\end{table}

Memory helps precisely where information is missing and nowhere else, ruling out increased capacity as the explanation.

\subsection{Negative result: landmark density in open worlds}

We tested whether the \emph{quantity} of persistent visual anchors quantitatively predicts abstraction quality, using the open-world environment with $n$ persistent landmarks among per-episode-reshuffled obstacles. In an earlier exploratory version of this pipeline, single runs at $n=0,3,6$ suggested an apparently clean monotonic trend. Replicating with three seeds per point did not support it (Table~\ref{tab:landmark}); the exploratory single runs are not part of the released results and we report no numbers from them.

\begin{table}[h]
\centering
\caption{Random-control gap by landmark count, mean $\pm$ std over 3 seeds. Source: \resfile{table3_landmarks.json}.}
\label{tab:landmark}
\begin{tabular}{lccccc}
\toprule
\tablerows{generated/table3.tex}
\bottomrule
\end{tabular}
\end{table}

No monotonic relationship survives; between-point differences are comparable to within-point standard deviation. We report this because the single-run version was convincing enough that we had begun writing it up as a positive finding.

What \emph{does} hold qualitatively across every variant tested: pure dead reckoning (zero landmarks, obstacles reshuffled every episode) never produces a large gap, while the building environment's doorways --- frequent, reliable, visually distinct anchors encountered on every room transition --- produce the largest gap observed ($\TEgoRecGap$pp, Table~\ref{tab:2x2}). We therefore conjecture, without claiming to have established, that anchor \emph{reliability and frequency of encounter} matter more than raw count.

\section{Discussion}

Our results suggest hierarchical abstraction becomes measurably real only when two conditions hold simultaneously: \textbf{a mechanism capable of integrating history}, and \textbf{information genuinely unavailable to the lower level without it}. Neither alone suffices, and each failure mode is invisible to standard metrics --- the full-observability cells post the highest absolute probe scores in the entire study while learning nothing.

This offers an empirical account of an architectural choice usually stated without justification: RSSM-style world models (e.g.\ DreamerV3, \citealp{hafner2023dreamerv3}) pair a recurrent deterministic state with an instantaneous stochastic latent. Our $2\times2$ indicates that on fully-observable benchmarks the recurrent component is approximately inert, and its contribution appears only under partial observability --- which is consistent with, and may partly explain, why hierarchical or recurrent components sometimes appear to contribute little on toy fully-observable tasks.

\textbf{Relation to end-to-end hierarchical JEPAs.} The environments of \citet{zhang2026hjepa} use a fixed camera, and their level encoders are pointwise functions of the level below; temporal context enters only through the predictors. In our terms this is the full-observability, instantaneous-abstractor cell, where a higher level can reshape or discard information but has nothing to add. This is consistent with their own attribution of planning gains to cost geometry and temporal decomposition rather than to information unavailable at level 1, and with their finding that agent position stays equally decodable at every level. Our results predict where a further, informational benefit of hierarchy should appear: under partial observability, and only if upper-level encoders integrate history. Applying the random-abstractor control to their trained levels is a direct test we leave to future work.

\subsection{Limitations}

All environments are small synthetic gridworlds with hand-specified slow variables, far simpler than the visual navigation and manipulation suites of \citet{zhang2026hjepa}; nothing here has touched real sensor data or learned its own abstraction hierarchy rather than being probed against a predefined one. The open-world ``slow variable'' is an artificial zone grid imposed for probing, not an emergent structure. The landmark-density question remains open rather than answered negatively --- our null result bounds the effect size detectable at our sample size, it does not demonstrate absence. All results are averaged over only three seeds; the seed-to-seed spread of the central gap is wide, and effect sizes should be read accordingly.

\section{Conclusion}

We provide a diagnostic protocol for verifying that hierarchical world-model abstractions are learned rather than decorative, and apply it to show that three of four natural design conditions produce abstractions indistinguishable from random projections. We argue the random-abstractor control is cheap enough --- one extra untrained module and one extra probe --- that it should be standard in any work claiming a hierarchical model has learned a meaningful abstraction.

Code, environments, all diagnostics, and the result files behind every number in this paper are available at \url{\https://github.com/r3gnnm/hier-abstraction-diagnostic}. A single command (\texttt{reproduce.sh} or \texttt{reproduce.bat}) regenerates all results and tables.

\appendix
\section{Reproduction}

\begin{verbatim}
python env_stats.py --episodes 500 --ep-len 75
python run_2x2_sweep.py --seeds 0 1 2 --episodes 200 --ep-len 48 --epochs 40 --k 8
python compare_memory.py --env egocentric --episodes 150 --ep-len 32 --epochs 20 --seeds 0 1 2
python compare_memory.py --env base       --episodes 150 --ep-len 32 --epochs 20 --seeds 0 1 2
python run_landmark_sweep.py --landmarks 0 3 6 10 15 --seeds 0 1 2 \
    --episodes 200 --ep-len 48 --epochs 40 --k 8
python make_tables.py      # results/*.json -> paper/generated/*.tex
\end{verbatim}

\bibliographystyle{plainnat}
\begin{thebibliography}{9}

\bibitem[Assran et~al.(2023)]{assran2023ijepa}
Assran, M., Duval, Q., Misra, I., Bojanowski, P., Vincent, P., Rabbat, M., LeCun, Y., and Ballas, N.
\newblock Self-supervised learning from images with a joint-embedding predictive architecture.
\newblock In \emph{Proceedings of the IEEE/CVF Conference on Computer Vision and Pattern Recognition (CVPR)}, pp.\ 15619--15629, 2023.

\bibitem[Bardes et~al.(2021)]{bardes2021vicreg}
Bardes, A., Ponce, J., and LeCun, Y.
\newblock VICReg: Variance-invariance-covariance regularization for self-supervised learning.
\newblock \emph{arXiv preprint arXiv:2105.04906}, 2021.

\bibitem[Bardes et~al.(2024)]{bardes2024vjepa}
Bardes, A., Garrido, Q., Ponce, J., Chen, X., Rabbat, M., LeCun, Y., Assran, M., and Ballas, N.
\newblock Revisiting feature prediction for learning visual representations from video.
\newblock \emph{arXiv preprint arXiv:2404.08471}, 2024.

\bibitem[Hafner et~al.(2023)]{hafner2023dreamerv3}
Hafner, D., Pasukonis, J., Ba, J., and Lillicrap, T.
\newblock Mastering diverse domains through world models.
\newblock \emph{arXiv preprint arXiv:2301.04104}, 2023.

\bibitem[LeCun(2022)]{lecun2022path}
LeCun, Y.
\newblock A path towards autonomous machine intelligence, version 0.9.2.
\newblock \emph{OpenReview}, 2022.
\newblock URL \url{https://openreview.net/pdf?id=BZ5a1r-kVsf}.

\bibitem[Zhang et~al.(2026)]{zhang2026hjepa}
Zhang, W., Terver, B., Rabbat, M., LeCun, Y., and Balestriero, R.
\newblock {H-JEPA}: End-to-end learning of hierarchical world models for visual planning.
\newblock \emph{arXiv preprint arXiv:2610.06805}, 2026.

\end{thebibliography}

\end{document}
